%% AUTO-GENERATED by doc/manuals/mcstas/gen_mcrun_options_table.py
%% from tools/Python/mcrun/mcrun.py -- DO NOT EDIT BY HAND.
%% Regenerate with:
%%   python3 gen_mcrun_options_table.py ../../../tools/Python/mcrun/mcrun.py > mcrun_options_table.tex

\begin{longtable}{|p{0.32\textwidth}|p{0.58\textwidth}|}
\caption{\texttt{mcrun}-specific options (compilation, scanning, optimisation, MPI/OpenACC).} \label{f:mcrunoptions} \\
\hline
\textbf{Option} & \textbf{Description} \\ \hline
\endfirsthead
\hline \textbf{Option} & \textbf{Description} \\ \hline
\endhead
\texttt{-c} \newline \texttt{--force-compile} & force rebuilding of instrument \\ \hline
\texttt{--cogen}~\textit{cogen} & Choice of code-generator (implies -c) \\ \hline
\texttt{-C} \newline \texttt{--c-lint} & Use c-linter (e.g. cppcheck) to lint the generated code. Configure linter via mccode\_config.json. Implies -c and -v, but also NO simulation will be run. \\ \hline
\texttt{-I} & Append to McCode search path (implies -c) \\ \hline
\texttt{--D1}~\textit{D1} & Set extra -D args (implies -c) \\ \hline
\texttt{--D2}~\textit{D2} & Set extra -D args (implies -c) \\ \hline
\texttt{--D3}~\textit{D3} & Set extra -D args (implies -c) \\ \hline
\texttt{-p} \newline \texttt{--param}~\textit{FILE} & Read parameters from file FILE \\ \hline
\texttt{-N} \newline \texttt{--numpoints}~\textit{NP} & Set number of scan points \\ \hline
\texttt{--seeds}~\textit{SEEDS} & Set range of seeds to scan (each must be: SEED != 0) \\ \hline
\texttt{-L} \newline \texttt{--list} & Use a fixed list of points for linear scanning \\ \hline
\texttt{-M} \newline \texttt{--multi} & Run a multi-dimensional scan \\ \hline
\texttt{--scan\_split}~\textit{scan\_split} & Scan by parallelising steps as individual cpu threads. Initialise by number of wanted threads (e.g. your number of cores). \\ \hline
\texttt{--autoplot} & Open plotter on generated dataset \\ \hline
\texttt{--invcanvas} & Forward request for inverted canvas to plotter \\ \hline
\texttt{--autoplotter} & Specify the plotter used with --autoplot \\ \hline
\texttt{--embed} & Store copy of instrument file in output directory \textit{(default: True)} \\ \hline
\texttt{--mpi}~\textit{NB\_CPU} & Spread simulation over NB\_CPU machines using MPI \\ \hline
\texttt{--openacc} & parallelize using openacc \\ \hline
\texttt{--funnel} & funneling simulation flow, e.g. for mixed CPU/GPU \\ \hline
\texttt{--machines}~\textit{machines} & Defines path of MPI machinefile to use in parallel mode \\ \hline
\texttt{--optimise-file}~\textit{FILE} & Store scan results in FILE (defaults to: "mccode.dat") \\ \hline
\texttt{--no-cflags} & Disable optimising compiler flags for faster compilation \\ \hline
\texttt{--no-main} & Do not generate a main(), e.g. for use with mcstas2vitess.pl. Implies -c \\ \hline
\texttt{--verbose} & Enable verbose output \\ \hline
\texttt{--write-user-config} & Generate a user config file \\ \hline
\texttt{--edit-user-config} & Generate and edit user config file in EDITOR \\ \hline
\texttt{--override-config}~\textit{PATH} & Load config file from specific dir \\ \hline
\texttt{--optimize} & Optimize instrument variable parameters to maximize monitors \\ \hline
\texttt{--optimize-maxiter}~\textit{optimize\_maxiter} & Maximum number of optimization iterations to perform. Default=1000 \textit{(default: 1000)} \\ \hline
\texttt{--optimize-tol}~\textit{optimize\_tol} & Tolerance for optimization termination. When optimize-tol is specified, the selected optimization algorithm sets some relevant solver-specific tolerance(s) equal to optimize-tol \\ \hline
\texttt{--optimize-method}~\textit{optimize\_method} & Optimization solver in [...] (default: [...]) You can use your custom method method(fun, x0, args, **kwargs, **options). Please refer to scipy documentation for proper use of it: https://docs.scipy.org/doc/scipy/reference/generated/scipy.optimize.minimize.html?highlight=minimize \\ \hline
\texttt{--optimize-eval}~\textit{optimize\_eval} & Optimization expression to evaluate for each detector "d" structure. You may combine: "d.intensity" The detector intensity; "d.error" The detector intensity uncertainty; "d.values" An array with [intensity, error, counts]; "d.X0 d.Y0" Center of signal (1st moment); "d.dX d.dY" Width of signal (2nd moment). Default is "d.intensity". Examples are: "d.intensity/d.dX" and "d.intensity/d.dX/d.dY" \\ \hline
\texttt{--optimize-minimize} & Choose to minimize the monitors instead of maximize \\ \hline
\texttt{--optimize-monitor}~\textit{optimize\_monitor} & Name of a single monitor to optimize (default is to use all) \\ \hline
\texttt{--showcfg}~\textit{ITEM} & Print selected cfg item and exit (paths are resolved and absolute). Allowed values are particle. \\ \hline
\end{longtable}



\begin{longtable}{|p{0.32\textwidth}|p{0.58\textwidth}|}
\caption{Instrument/simulation options (also accepted directly by the generated \texttt{<instr>.out} executable).} \label{f:simoptions} \\
\hline
\textbf{Option} & \textbf{Description} \\ \hline
\endfirsthead
\hline \textbf{Option} & \textbf{Description} \\ \hline
\endhead
\texttt{-s} \newline \texttt{--seed}~\textit{SEED} & Set random seed (must be: SEED != 0) \\ \hline
\texttt{-n} \newline \texttt{--ncount}~\textit{COUNT} & Set number of particles to simulate \textit{(default: 1000000)} \\ \hline
\texttt{-t} \newline \texttt{--trace}~\textit{trace} & Enable trace of particles through instrument \textit{(default: 0)} \\ \hline
\texttt{--no-trace}~\textit{notrace} & Disable trace of particles in instrument (combine with -c) \\ \hline
\texttt{-y} \newline \texttt{--yes} & Assume any default parameter value in instrument \\ \hline
\texttt{-d} \newline \texttt{--dir}~\textit{DIR} & Put all data files in directory DIR. If unspecified INSTRUMENT\_TIMESTAMP is used \\ \hline
\texttt{--dirprefix}~\textit{dirprefix} & Put all data files in directory PREFIX\_TIMESTAMP \\ \hline
\texttt{--dirsuffix}~\textit{dirsuffix} & Put all data files in directory INSTRUMENT\_DIRSUFFIX \\ \hline
\texttt{-a} \newline \texttt{--append} & Append data files to those already in directory DIR \\ \hline
\texttt{--format}~\textit{FORMAT} & Output data files using format FORMAT, usually McCode or NeXus (format list obtained from <instr>.out -h) \textit{(default: McCode)} \\ \hline
\texttt{--bufsiz}~\textit{BUFSIZ} & Monitor\_nD list/buffer-size (defaults to [...]) \\ \hline
\texttt{--vecsize}~\textit{VECSIZE} & vector length in OpenACC parallel scenarios \\ \hline
\texttt{--numgangs}~\textit{NUMGANGS} & number of 'gangs' in OpenACC parallel scenarios \\ \hline
\texttt{--gpu\_innerloop}~\textit{INNERLOOP} & Maximum particles in an OpenACC kernel run. (If INNERLOOP is smaller than ncount we repeat) \\ \hline
\texttt{--no-output-files} & Do not write any data files \\ \hline
\texttt{-i} \newline \texttt{--info} & Detailed instrument information \\ \hline
\texttt{--list-parameters} & Print the instrument parameters to standard out \\ \hline
\texttt{--meta-list} & Print all metadata defining component names \\ \hline
\texttt{--meta-defined} & Print metadata names for component, or indicate if component:name exists \\ \hline
\texttt{--meta-type} & Print metadata type for component:name \\ \hline
\texttt{--meta-data} & Print metadata for component:name \\ \hline
\texttt{-g} \newline \texttt{--gravitation} \newline \texttt{--gravity} & Enable gravitation for all trajectories \\ \hline
\texttt{--IDF} & Flag to attempt inclusion of XML-based IDF when --format=NeXus (format list obtained from <instr>.out -h) \\ \hline
\end{longtable}
